\documentclass[11pt]{article}
\usepackage[margin=25mm]{geometry}
\usepackage{fontspec}
\setmainfont{Latin Modern Roman}
\usepackage{amsmath,amssymb,mathtools}
\usepackage{unicode-math}
\setmathfont{Latin Modern Math}
\usepackage{xurl,hyperref}
\hypersetup{colorlinks=true,urlcolor=blue}
\setcounter{secnumdepth}{0}
\setlength{\parindent}{0pt}
\setlength{\parskip}{5pt}
\setlength{\emergencystretch}{3em}
\newcommand{\sref}[2]{\hyperlink{source:#1:#2}{[S#2]}}
\newcommand{\cataloguescope}[1]{\paragraph{Recorded target.}#1}
\begin{document}
\section{Hadamard Conjecture}\label{statement}

\subsection{Definitions and mathematical statement}

\paragraph{Shared notation.}$\mathbb N=\{1,2,\ldots\}$, and $\mathbb Z,\mathbb Q,\mathbb R,\mathbb C$ denote the integers, rationals, reals and complex numbers. Any other conventions are specified locally.


A Hadamard matrix of order $m$ is a matrix $H=(h_{ij})\in\{-1,1\}^{m\times m}$ satisfying $HH^{\mathsf T}=mI_m$, where $H^{\mathsf T}$ is its transpose and $I_m$ is the identity matrix. Equivalently, its rows are pairwise orthogonal for the usual real dot product; each row automatically has squared length $m$.
\paragraph{Statement recorded in the supplied source.}
For every integer $k\geq1$, does there exist a Hadamard matrix $H$ of order $4k$? Thus the required equality is $HH^{\mathsf T}=4kI_{4k}$ over the integers, rather than a congruence modulo a fixed modulus. \sref{statement}{2}
\subsection{Short English statement}
Does every positive order divisible by four admit a square sign matrix with mutually orthogonal rows?
\subsection{Sources}
\begin{description}
\item[\hypertarget{source:statement:1}{S1}] ULAM.AI and the UnsolvedMath Contributors, \emph{UnsolvedMath: A Curated Collection of Open Mathematics Problems}, catalogue entry. CC BY 4.0. \url{https://huggingface.co/datasets/ulamai/UnsolvedMath}.
\item[\hypertarget{source:statement:2}{S2}] Shalom Eliahou, \emph{A 64-modular Hadamard matrix of order 668}, Australasian Journal of Combinatorics 93(2) (2025), 422--427, §1, p.422. \url{https://ajc.maths.uq.edu.au/pdf/93/ajc_v93_p422.pdf}.
\end{description}
\label{end:statement}
\end{document}
